\section{Study 2: Is the reliability landscape a moving target?}\label{sec:s1b}
% =====================================================================
A systematic review necessarily captures results \emph{as reported}, and most ``promising''-band evidence in
Study~1 was produced with proprietary judges that are now one or two model generations old (e.g.\ the influential
G-Eval used GPT-4). This raises a question that the review alone cannot answer: when an LLM judge falls short of the
human ceiling in a promising domain, is that a \emph{temporary} limitation that more capable models will erase
(``matter of time''), or a \emph{structural} property of LLM-as-a-judge that better models do not overcome? Study~2
answers this directly for a flagship case by re-running the evaluation with the most powerful open models available
to us today and comparing against the canonical reported numbers.

\subsection{Method}
We select \textbf{SummEval} \citep{fabbri2021summeval} (summarization quality assessment, a central amber-band task
in Study~1) because it pairs (i) a widely cited LLM-judge result to beat (G-Eval/GPT-4; \citealp{liu2023geval})
with (ii) public expert human ratings on four aspects (coherence, consistency, fluency, relevance) for
$100\times16=1600$ machine summaries. We replicate the G-Eval summary-level protocol \emph{verbatim} (same aspect
definitions, 1-5 integer ratings, temperature~0) but swap the judge for the \ONEmodels{} most powerful open models
in the available open-model catalogue (\ONEpanel). The comparison metric is the canonical \emph{summary-level Spearman}: for each
source document, correlate the judge's 16 scores with the human means, averaged across documents. The bar to beat is
G-Eval/GPT-4's reported per-aspect correlations (coherence~.582, consistency~.507, fluency~.455, relevance~.547;
average~$\approx$.52). In total Study~2 contributes \ONEnjudg{} fresh judgments.

\subsection{Results}
\begin{table}[t]\centering\small
\caption{Study 2. Summary-level Spearman $\rho$ of the most powerful open judges on SummEval vs.\ the reported
G-Eval-4/GPT-4 correlations (\ONEnjudg{} judgments). Green cells \emph{surpass} the reported per-aspect number.}
\label{tab:study1b}
\begin{tabular}{@{}l cccc c@{}}
\toprule
Judge & Coher. & Consist. & Fluency & Relev. & Avg $\rho$ \\
\midrule
\textit{Reported (G-Eval-4)} & 0.582 & 0.507 & 0.455 & 0.547 & 0.523 \\
\midrule
\texttt{deepseek-v4-pro} & 0.548 & \cellcolor{g!30}0.688 & 0.392 & 0.540 & \cellcolor{g!30}0.542 \\
\texttt{qwen3.6-35b-a3b} & \cellcolor{g!30}0.606 & \cellcolor{g!30}0.702 & 0.298 & 0.540 & \cellcolor{g!30}0.536 \\
\texttt{glm-5.1} & \cellcolor{g!30}0.650 & \cellcolor{g!30}0.693 & 0.217 & 0.543 & \cellcolor{g!30}0.526 \\
\texttt{gptoss-120b} & \cellcolor{g!30}0.586 & \cellcolor{g!30}0.685 & 0.260 & 0.482 & 0.503 \\
\bottomrule
\end{tabular}
\end{table}
We follow the standard summary-level G-Eval protocol \citep{liu2023geval} and attach bootstrap 95\% CIs (over the
100 source documents) to every correlation (Figure~\ref{fig:s1b}, Table~\ref{tab:study1b}). On the verifiable
aspects the open judges \emph{meet or exceed} the established G-Eval/GPT-4 benchmark: \RTwobeats{} of the 16
model$\times$aspect cells have a 95\% CI lying \emph{entirely above} the reported value, all on \textbf{consistency}
and \textbf{coherence} (e.g.\ \ONEbest{} consistency $\rho{=}0.69$, CI $[0.63,0.74]$ vs.\ $.507$), and
\textbf{relevance} is on par ($.48$-$.54$ vs.\ $.55$). Our protocol differs from the original in two deliberate,
deployment-relevant ways: the judges are open and locally-servable, and we read a single integer score rather than a
proprietary model's probability-weighted decoding. The comparison is to the figures reported by the original work,
the standard reference point in this literature. But on \textbf{fluency} every open model falls well
\emph{short} ($.22$-$.39$ vs.\ a reported $.455$). We diagnose why: fluency is the most \emph{compressed} aspect of
the human reference, with the lowest within-document rating variance ($\mathrm{SD}=\SEfluSD$ vs $\SEconSD$ for
consistency) and $\SEfluTop\%$ of summaries rated near the top of the scale, so it offers the least to rank, and it
is the aspect on which the \emph{reported proprietary judge also scores lowest}. This is a property of the
\emph{reference}, not the judges: even a noiseless judge returning the rounded human means tops out at
$\rho{\approx}\SEfluSelf$ on fluency. The clean result is therefore that capability gains concentrate on the
verifiable, rankable aspects (consistency, coherence), while the most compressed aspect is hard to rank for every
judge, open or proprietary.

\begin{figure}[t]\centering
\includegraphics[width=0.97\linewidth]{s1b_fig01_spearman}
\caption{Study 2. Summary-level Spearman of the most powerful open judges on SummEval, per aspect, against the
reported G-Eval/GPT-4 bar (dashed). Bars above the dashed line indicate an open model that \emph{surpasses} the
reported proprietary-judge correlation.}
\label{fig:s1b}
\end{figure}

\paragraph{Interpretation: time vs.\ structure.} Study~2 shows the reliability landscape is \emph{partly} a moving
target. Where Study~1's amber rating reflected a capability gap on a verifiable sub-judgment, today's open models
close or exceed it; reliability there is largely a matter of time and access. But where the shortfall reflects a
\emph{compressed or contested human ceiling} (e.g.\ fluency, and, as Study~3 will show, subjective social annotation),
even the most powerful models do not break through, because the limit lives in the human reference, not the model.
This distinction, temporal capability gap vs.\ structural ceiling, is the lens we carry into HCC.

\paragraph{Takeaway.} Study~1 says LLM judges are safe where the task is objective and verifiable, and
Study~2 shows that frontier open models are now strong enough to match proprietary judges on exactly those
verifiable tasks. But subjective, high-disagreement constructs (the kind that
humans themselves contest) are precisely the tasks Study~1 flags as risky and Study~2
predicts scale will \emph{not} fix. Study~3 then tests this prediction directly, on datasets with rich human label distributions.

% =====================================================================
